\documentclass[11pt,a4paper]{article}

% Encoding and language
\usepackage[utf8]{inputenc}
\usepackage[T5]{fontenc}      % for Vietnamese
\usepackage[vietnamese,english]{babel}  % choose as needed

% Margins and layout
\usepackage[top=2cm,bottom=2cm,left=2.5cm,right=2.5cm]{geometry}
\usepackage{parskip}          % No indent, space between paragraphs
\usepackage{multicol}         % Optional for two-column layout
\usepackage{lastpage}

% Fonts
\usepackage{lmodern}          % Clean scalable font
\usepackage{microtype}        % Better typography
\usepackage{titlesec}         % Custom section formatting
\usepackage[normalem]{ulem} % <- strikethrough

\usepackage{tikz}
\usetikzlibrary{arrows}
\usetikzlibrary{arrows.meta, positioning}
\usetikzlibrary{calc}
\usetikzlibrary{patterns}

% Code and input/output formatting
\usepackage{listings}
\usepackage{xcolor}
\lstset{
    basicstyle=\ttfamily\small,
    backgroundcolor=\color{gray!10},
    frame=single,
    columns=fullflexible,
    keepspaces=true,
    showstringspaces=false,
    breaklines=true,
    language=C++,
    tabsize=4,
}

% Math
\usepackage{amsmath, amssymb}
\usepackage{mathtools}

% Links
\usepackage[hyperfootnotes=false]{hyperref}

% Boxes and emphasis
\usepackage{tcolorbox}
\tcbuselibrary{listings, breakable}
\tcbset{colback=gray!5, colframe=black!20, boxrule=0.5pt}

% Title formatting
\titleformat{\section}{\large\bfseries}{\thesection}{1em}{}
\titleformat{\subsection}{\normalsize\bfseries}{\thesubsection}{1em}{}

% Misc
\usepackage{array}

% Header/footer (optional)
\usepackage{fancyhdr}
\pagestyle{fancy}
\fancyhf{}
\lhead{Library Checker -- Phan Thành Hưng -- Trường Phổ thông Năng khiếu, ĐHQG TP.HCM}
\renewcommand{\headrulewidth}{0.4pt}
\renewcommand{\footrulewidth}{0.4pt}

\usepackage{xifthen}  % dùng \ifthenelse và \isempty

% Macro ví dụ với số và giải thích đẹp
\newcommand{\ExampleBox}[4]{%
    \vspace{0.5em}
    \begin{tcolorbox}[title=\textcolor{darkgray}{Ví dụ~#1}, breakable]
    \ttfamily
    \textbf{Input} \\
    #2

    \vspace{0.8em}
    \textbf{Output} \\
    #3

    \ifthenelse{\isempty{#4}}{}{
        \vspace{0.8em}
        \textbf{Giải thích} \\
        \normalfont #4
    }
    \end{tcolorbox}
}

% PROBLEM CREATER TEMPLATE
\newcommand{\makeproblem}[2]{%
  \section{#1\ - #2}%
}

\newcommand{\inputDesc}[1]{%
  \vspace{0.8em}
  \noindent{\sffamily \textbf{Dữ liệu vào. }}%
  \ifthenelse{\isempty{#1}}%
    {}%
    {Đọc từ tệp \texttt{#1.INP}:}%
}

\newcommand{\outputDesc}[1]{%
  \vspace{0.8em}
  \noindent{\sffamily \textbf{Dữ liệu ra. }}%
  \ifthenelse{\isempty{#1}}%
    {}%
    {Xuất ra tệp \texttt{#1.OUT}:}%
}
%=============================

\title{LIBRARY CHECKER}
\author{Tác giả: Phan Thành Hưng}

\begin{document}
\maketitle
% %==========================

\vfill
\begin{center}
    \small
    \textit{Tài liệu gồm \ref{LastProblem} bài, có \pageref{LastPage} trang.}
\end{center}
\vspace{5em}

%==========================

\newpage
\def\ProblemCode{SEGTREE}
\def\ProblemTitle{Segment Tree (Basic)}
\makeproblem{\ProblemCode}{\ProblemTitle}

Cho dãy số nguyên $a_1, a_2, \dots, a_n$.
Xử lý $q$ truy vấn thuộc một trong các dạng sau:

\begin{itemize}
    \item \boxed{1\ k\ x}: Gán phần tử thứ $k$ thành $x\ (|x| \le 10^{9})$.
    \item \boxed{2\ l\ r}: Tính tổng các phần tử trong đoạn $[l, r]$.
    \item \boxed{3\ l\ r}: Tìm phần tử có giá trị nhỏ nhất trong đoạn $[l, r]$.
    \item \boxed{4\ l\ r}: Tìm phần tử có giá trị lớn nhất trong đoạn $[l, r]$.
\end{itemize}

% ===== INPUT DESCRIPTION =====
\inputDesc{}

\begin{itemize}
    \item Dòng đầu tiên gồm hai số nguyên dương $n, q\ (1 \le n, q \le 10^{5})$.
    \item Dòng thứ hai gồm $n$ số nguyên $a_1, a_2, \dots, a_n\ (|a_i| \le 10^{9})$.
    \item $q$ dòng tiếp theo, mỗi dòng chứa một truy vấn.
\end{itemize}

% ===== OUTPUT DESCRIPTION =====
\outputDesc{}

\begin{itemize}
    \item Với mỗi truy vấn loại 2, 3 hoặc 4, in ra một dòng chứa kết quả tương ứng.
\end{itemize}

% ===== SUBTASKS =====
% \noindent{\sffamily \textbf{Subtasks}}
% \begin{itemize}
%     \item \textbf{Subtask ?} ($?\%$ số điểm): Không có ràng buộc gì thêm.
% \end{itemize}

\vspace{0.5em}

\ExampleBox{1}{
8 10\\
5 -1 3 7 -4 2 6 -8\\
2 1 8\\
3 2 5\\
4 4 7\\
1 5 10\\
2 3 6\\
1 8 4\\
2 6 8\\
3 1 4\\
1 2 -9\\
4 1 3
}{
10\\
-4\\
7\\
22\\
12\\
-1\\
5
}{}


\newpage
\def\ProblemCode{SEGLAZY}
\def\ProblemTitle{Segment Tree (Lazy propagation)}
\makeproblem{\ProblemCode}{\ProblemTitle}

Cho dãy số nguyên $a_1, a_2, \dots, a_n$.
Xử lý $q$ truy vấn thuộc một trong các dạng sau:

\begin{itemize}
    \item \boxed{1\ l\ r\ k}: Tăng mỗi phần tử trong đoạn $[l, r]$ lên $k$ đơn vị $(|k| \le 10^{9})$.
    \item \boxed{2\ l\ r}: Tính tổng các phần tử trong đoạn $[l, r]$.
    \item \boxed{3\ l\ r}: Tìm phần tử có giá trị nhỏ nhất trong đoạn $[l, r]$.
    \item \boxed{4\ l\ r}: Tìm phần tử có giá trị lớn nhất trong đoạn $[l, r]$.
\end{itemize}

% ===== INPUT DESCRIPTION =====
\inputDesc{}

\begin{itemize}
    \item Dòng đầu tiên gồm hai số nguyên dương $n, q\ (1 \le n, q \le 10^{5})$.
    \item Dòng thứ hai gồm $n$ số nguyên $a_1, a_2, \dots, a_n\ (|a_i| \le 10^{9})$.
    \item $q$ dòng tiếp theo, mỗi dòng chứa một truy vấn.
\end{itemize}

% ===== OUTPUT DESCRIPTION =====
\outputDesc{}

\begin{itemize}
    \item Với mỗi truy vấn loại 2, 3 hoặc 4, in ra một dòng chứa kết quả tương ứng.
\end{itemize}

% ===== SUBTASKS =====
% \noindent{\sffamily \textbf{Subtasks}}
% \begin{itemize}
%     \item \textbf{Subtask ?} ($?\%$ số điểm): Không có ràng buộc gì thêm.
% \end{itemize}

\vspace{0.5em}

\ExampleBox{1}{
5 9\\
-3 5 -2 4 -1\\
2 1 5\\
3 1 5\\
4 1 5\\
1 2 4 3\\
2 1 5\\
3 2 4\\
4 2 4\\
1 1 5 -2\\
2 1 3
}{
3\\
-3\\
5\\
12\\
1\\
8\\
0
}{}


\newpage
\def\ProblemCode{QUERY2D}
\def\ProblemTitle{2D Data Structure (Basic)}
\makeproblem{\ProblemCode}{\ProblemTitle}

Cho lưới ô vuông kích thước $n \times m$. Ô nằm tại hàng $i$, cột $j$ có giá trị là $a_{i,j}$.

Xử lý $q$ truy vấn thuộc một trong các dạng sau:

\begin{itemize}
    \item \boxed{1\ x\ y\ k}: Tăng ô có tọa độ $(x,y)$ lên $k$ đơn vị $(|k| \le 10^{9})$.
    \item \boxed{2\ x_1\ y_1\ x_2\ y_2}: Tính tổng giá trị các ô vuông trong hình chữ nhật từ $(x_1,y_1)$ tới $(x_2,y_2)$.
\end{itemize}

% ===== INPUT DESCRIPTION =====
\inputDesc{}

\begin{itemize}
    \item Dòng đầu tiên gồm hai số nguyên dương $n, m, q\ (1 \le n, m \le 10^{3},\ 1\le q \le 10^{5})$.
    \item Dòng thứ $i$ trong $n$ dòng tiếp theo chứa $m$ số nguyên $a_{i,1}, a_{i,2}, \dots, a_{i,n}\ (|a_{i,j}| \le 10^{9})$.
    \item $q$ dòng tiếp theo, mỗi dòng chứa một truy vấn.
\end{itemize}

% ===== OUTPUT DESCRIPTION =====
\outputDesc{}

\begin{itemize}
    \item Với mỗi truy vấn loại 2, in ra một dòng chứa kết quả tương ứng.
\end{itemize}

% ===== SUBTASKS =====
% \noindent{\sffamily \textbf{Subtasks}}
% \begin{itemize}
%     \item \textbf{Subtask ?} ($?\%$ số điểm): Không có ràng buộc gì thêm.
% \end{itemize}

\vspace{0.5em}

\ExampleBox{1}{
4 4 7\\
1 2 3 4\\
5 6 7 8\\
9 10 11 12\\
13 14 15 16\\
2 1 1 4 4\\
1 2 2 5\\
2 2 2 3 3\\
1 4 4 -10\\
2 3 3 4 4\\
1 1 1 3\\
2 1 1 2 2
}{
136\\
39\\
44\\
22
}{}


\newpage
\def\ProblemCode{SPARSEBIT2D}
\def\ProblemTitle{Fenwick Tree 2D (Compressed)}
\makeproblem{\ProblemCode}{\ProblemTitle}

Cho một lưới ô vuông vô hạn trên mặt phẳng tọa độ. Ban đầu, giá trị tại mọi ô đều bằng $0$.

Bạn cần xử lý $q$ truy vấn, mỗi truy vấn thuộc một trong hai loại sau:

\begin{itemize}
    \item \boxed{1\ x\ y\ k}: Tăng ô có tọa độ $(x,y)$ lên $k$ đơn vị $(1 \le x, y \le 10^9, |k| \le 10^{9})$.
    \item \boxed{2\ x_1\ y_1\ x_2\ y_2}: Tính tổng giá trị các ô vuông trong hình chữ nhật từ $(x_1,y_1)$ tới $(x_2,y_2)$ $(1 \le x_1 \le x_2 \le 10^9,\ 1 \le y_1 \le y_2 \le 10^9)$.
\end{itemize}

% ===== INPUT DESCRIPTION =====
\inputDesc{}

\begin{itemize}
    \item Dòng đầu tiên gồm hai số nguyên dương $q\ (1\le q \le 2 \times 10^{5})$.
    \item $q$ dòng tiếp theo, mỗi dòng chứa một truy vấn.
\end{itemize}

% ===== OUTPUT DESCRIPTION =====
\outputDesc{}

\begin{itemize}
    \item Với mỗi truy vấn loại 2, in ra một dòng chứa kết quả tương ứng.
\end{itemize}

% ===== SUBTASKS =====
% \noindent{\sffamily \textbf{Subtasks}}
% \begin{itemize}
%     \item \textbf{Subtask ?} ($?\%$ số điểm): Không có ràng buộc gì thêm.
% \end{itemize}

\vspace{0.5em}

\ExampleBox{1}{
6\\
1 1 1 5\\
1 2 3 7\\
2 1 1 3 4\\
1 2 2 4\\
2 2 2 3 3\\
2 1 1 5 5
}{
12\\
11\\
16
}{}


\newpage
\input{SEGPERS}

\newpage
\def\ProblemCode{SEGPERSLAZY}
\def\ProblemTitle{Lazy Persistent Segment Tree}
\makeproblem{\ProblemCode}{\ProblemTitle}

Cho dãy số nguyên $a_1, a_2, \dots, a_n$.
Ban đầu, tồn tại phiên bản $1$ của dãy, chính là dãy ban đầu.
Mỗi truy vấn cập nhật sẽ tạo ra một phiên bản mới của dãy.

Xử lý $q$ truy vấn thuộc một trong các dạng sau:

\begin{itemize}
    \item \boxed{1\ v\ l\ r\ k}:
    Từ phiên bản $v$, tạo ra một phiên bản mới bằng cách tăng mỗi phần tử
    trong đoạn $[l, r]$ lên $k$ đơn vị $(|k| \le 10^{9})$.

    \item \boxed{2\ v\ l\ r}:
    Tính tổng các phần tử trong đoạn $[l, r]$ của phiên bản $v$.

    \item \boxed{3\ v\ l\ r}:
    Tìm giá trị nhỏ nhất trong đoạn $[l, r]$ của phiên bản $v$.

    \item \boxed{4\ v\ l\ r}:
    Tìm giá trị lớn nhất trong đoạn $[l, r]$ của phiên bản $v$.
\end{itemize}

Các phiên bản được đánh số liên tiếp từ $1$ theo thứ tự xuất hiện
(các truy vấn loại \texttt{1}).

% ===== INPUT DESCRIPTION =====
\inputDesc{}

\begin{itemize}
    \item Dòng đầu tiên gồm hai số nguyên dương $n, q\ (1 \le n, q \le 10^{5})$.
    \item Dòng thứ hai gồm $n$ số nguyên $a_1, a_2, \dots, a_n\ (|a_i| \le 10^{9})$.
    \item $q$ dòng tiếp theo, mỗi dòng chứa một truy vấn như mô tả ở trên.
\end{itemize}

% ===== OUTPUT DESCRIPTION =====
\outputDesc{}

\begin{itemize}
    \item Với mỗi truy vấn loại 2, 3 hoặc 4, in ra một dòng chứa kết quả tương ứng.
\end{itemize}

\vspace{0.5em}

\ExampleBox{2}{
6 9\\
2 -5 4 1 7 -3\\
2 1 1 6\\
1 1 2 5 3\\
2 2 1 6\\
3 2 3 6\\
4 2 1 4\\
1 2 4 6 -2\\
2 3 4 6\\
3 1 1 6\\
4 3 1 6
}{
6\\
18\\
-3\\
7\\
-3\\
-5\\
7
}{}


\newpage
\def\ProblemCode{WAVELET}
\def\ProblemTitle{Wavelet Tree}
\makeproblem{\ProblemCode}{\ProblemTitle}

Cho dãy số nguyên $a_1, a_2, \dots, a_n$. Hãy xử lý $q$ truy vấn thuộc một trong các dạng sau:

\begin{itemize}
    \item \boxed{1\ l\ r\ k}: Phần tử lớn thứ $k$ trong đoạn $[l, r]$.
    Nếu không tồn tại, in ra \texttt{NO}.
    \item \boxed{2\ l\ r\ x}: Số phần tử $\le x$ trong đoạn $[l, r]$.
    \item \boxed{3\ l\ r\ x}: Tổng các phần tử $\le x$ trong đoạn $[l, r]$.
    \item \boxed{4\ l\ r\ x}: Số phần tử $= x$ trong đoạn $[l, r]$.
\end{itemize}   

% ===== INPUT DESCRIPTION =====
\inputDesc{}

\begin{itemize}
    \item Dòng đầu tiên gồm hai số nguyên dương $n, q\ (1 \le n, q \le 10^{5})$.
    \item Dòng thứ hai gồm $n$ số nguyên $a_1, a_2, \dots, a_n\ (|a_i| \le 10^{6})$.
    \item $q$ dòng tiếp theo, mỗi dòng chứa một truy vấn.
\end{itemize}

% ===== OUTPUT DESCRIPTION =====
\outputDesc{}

\begin{itemize}
    \item Gồm $q$ dòng, mỗi dòng trả lời kết quả của truy vấn tương ứng.
\end{itemize}

% ===== SUBTASKS =====
% \noindent{\sffamily \textbf{Subtasks}}
% \begin{itemize}
%     \item \textbf{Subtask ?} ($?\%$ số điểm): Không có ràng buộc gì thêm.
% \end{itemize}

\vspace{0.5em}

\ExampleBox{1}{
6 8\\
-5 3 -2 7 -2 4\\
1 1 6 3\\
2 1 6 -2\\
3 1 6 0\\
4 1 6 -2\\
1 2 5 5\\
2 2 5 10\\
3 2 5 3\\
4 3 6 7
}{
3\\
3\\
-9\\
2\\
NO\\
4\\
-1\\
1
}{}


\newpage
\def\ProblemCode{WAVELET2}
\def\ProblemTitle{Wavelet Tree 2}
\makeproblem{\ProblemCode}{\ProblemTitle}

Cho một mảng $a$ gồm $n$ số nguyên.

Hãy xử lý $q$ truy vấn, mỗi truy vấn có dạng:

\begin{itemize}
    \item \boxed{l\ r\ k}:
    Tìm số nguyên $x$ nhỏ nhất sao cho $x$ xuất hiện trong đoạn $[l, r]$
    \textbf{nhiều hơn} $\dfrac{r-l+1}{k}$ lần.
    Nếu không tồn tại số như vậy, in ra $-1$.
\end{itemize}

% ===== INPUT DESCRIPTION =====
\inputDesc{}

\begin{itemize}
    \item Dòng đầu gồm hai số nguyên $n, q\ (1 \le n, q \le 3 \times 10^{5})$.
    \item Dòng thứ hai gồm $n$ số nguyên $a_1, a_2, \dots, a_n\ (1 \le a_i \le n)$.
    \item $q$ dòng tiếp theo, mỗi dòng gồm ba số nguyên
    $l, r, k\ (1 \le l \le r \le n,\ 2 \le k \le 5)$.
\end{itemize}

% ===== OUTPUT DESCRIPTION =====
\outputDesc{}

\begin{itemize}
    \item Với mỗi truy vấn, in ra một dòng chứa đáp án tương ứng.
\end{itemize}

\vspace{0.5em}

\ExampleBox{1}{
4 2\\
1 1 2 2\\
1 3 2\\
1 4 2
}{
1\\
-1
}{}

\ExampleBox{2}{
5 3\\
1 2 1 3 2\\
2 5 3\\
1 2 3\\
5 5 2
}{
2\\
1\\
2
}{}


\newpage
\def\ProblemCode{DSUROLLBACK}
\def\ProblemTitle{DSU Rollback (Basic)}
\makeproblem{\ProblemCode}{\ProblemTitle}

Cho đồ thị vô hướng gồm $n$ đỉnh. Các đỉnh được đánh số từ $1$ tới $n$. Ban đầu, đồ thị chưa có cạnh nào và các đỉnh $1,2,\dots,n$ lần lượt mang giá trị $a_1,a_2,\dots,a_n$.

Hãy xử lý $q$ truy vấn thuộc một trong các dạng sau:

\begin{itemize}
    \item \boxed{0\ u\ v}: Thêm một cạnh mới nối giữa hai đỉnh $u$ và $v$.
    \item \boxed{1}: Xóa cạnh gần nhất đã được thêm vào từ một thao tác loại \texttt{0} mà chưa bị xóa.
    \item \boxed{2}: Đếm số thành phần liên thông trong đồ thị.
    \item \boxed{3\ u}: Đếm số đỉnh thuộc cùng thành phần liên thông với đỉnh $u$ (bao gồm cả $u$).
    \item \boxed{4\ u}: Tìm giá trị đỉnh nhỏ nhất trong các đỉnh thuộc cùng thành phần liên thông với $u$.
\end{itemize}

\textbf{Lưu ý:} Giữa hai đỉnh có thể có nhiều cạnh nối phân biệt.

% ===== INPUT DESCRIPTION =====
\inputDesc{}

\begin{itemize}
    \item Dòng đầu tiên gồm hai số nguyên dương $n, q\ (1 \le n,q \le 10^{5})$.
    \item Dòng tiếp theo gồm $n$ số nguyên $a_1,a_2,\dots,a_n\ (|a_i| \le 10^{9})$.
    \item $q$ dòng tiếp theo, mỗi dòng chứa một truy vấn.
\end{itemize}

Dữ liệu đảm bảo tại mọi thời điểm, số truy vấn loại \texttt{0} xuất hiện không ít hơn số truy vấn loại \texttt{1}.

% ===== OUTPUT DESCRIPTION =====
\outputDesc{}

\begin{itemize}
    \item Với mỗi truy vấn loại 2, 3 hoặc 4, in ra một dòng chứa kết quả tương ứng.
\end{itemize}

% ===== SUBTASKS =====
% \noindent{\sffamily \textbf{Subtasks}}
% \begin{itemize}
%     \item \textbf{Subtask ?} ($?\%$ số điểm): Không có ràng buộc gì thêm.
% \end{itemize}

\vspace{0.5em}

\ExampleBox{1}{
5 8\\
3 -1 4 2 0\\
2\\
0 1 2\\
4 1\\
0 2 3\\
3 1\\ 
1\\
2\\
4 3
}{
5\\
-1\\
3\\
4\\
4
}{}


\newpage
\def\ProblemCode{VIRTREE}
\def\ProblemTitle{Virtual Tree (Basic)}
\makeproblem{\ProblemCode}{\ProblemTitle}

Cho cây vô hướng gồm $n$ đỉnh có gốc tại đỉnh $1$.

Hãy xử lý $q$ truy vấn. Mỗi truy vấn có dạng:

\begin{itemize}
    \item \boxed{k\ v_1\ v_2\ \dots\ v_k}:
    Cho $k$ đỉnh phân biệt trên cây. Xây dựng \textbf{Virtual Tree} chứa tất cả các đỉnh $v_i$ và các LCA cần thiết.
\end{itemize}

% ===== INPUT DESCRIPTION =====
\inputDesc{}

\begin{itemize}
    \item Dòng đầu tiên gồm hai số nguyên dương $n, q\ (1 \le n, q \le 2\times10^{5})$.
    \item $n-1$ dòng tiếp theo, mỗi dòng gồm hai số $u, v$ mô tả một cạnh của cây.
    \item $q$ dòng tiếp theo, mỗi dòng chứa một truy vấn.
\end{itemize}

Dữ liệu đảm bảo: $\sum k \le 5\times10^{5}$.

% ===== OUTPUT DESCRIPTION =====
\outputDesc{}

\begin{itemize}
    \item Với mỗi truy vấn, in ra:
    \begin{itemize}
        \item[+] Dòng đầu: một số $d$ là số đỉnh của Virtual Tree tương ứng.
        \item[+] $d-1$ dòng tiếp theo: mỗi dòng chứa một cạnh $u\ v$ trên Virtual Tree đó.
    \end{itemize}
    Bạn được phép in các cạnh theo thứ tự tùy ý.
\end{itemize}

% ===== SUBTASKS =====
% \noindent{\sffamily \textbf{Subtasks}}
% \begin{itemize}
%     \item \textbf{Subtask ?} ($?\%$ số điểm): Không có ràng buộc gì thêm.
% \end{itemize}

\vspace{0.5em}

\ExampleBox{1}{
7 2\\
1 2\\
1 3\\
2 4\\
2 5\\
3 6\\
3 7\\
3 4 5 6\\
2 4 7
}{
5\\
2 4\\
2 5\\
1 2\\
1 6\\
3\\
1 4\\
1 7
}{}


\newpage
\def\ProblemCode{DECENTROID}
\def\ProblemTitle{Centroid Decomposition}
\makeproblem{\ProblemCode}{\ProblemTitle}

Cho một cây gồm $n$ đỉnh được đánh số từ $1$ đến $n$. Hãy xây dựng cây phân rã trọng tâm (centroid decomposition) của cây đã cho.

% ===== INPUT DESCRIPTION =====
\inputDesc{}

\begin{itemize}
    \item Dòng đầu tiên gồm một số nguyên dương $n\ (1 \le n \le 2 \cdot 10^5)$.
    \item $n-1$ dòng tiếp theo, mỗi dòng gồm hai số nguyên $u, v\ (1 \le u, v \le n)$ mô tả một cạnh của cây.
\end{itemize}

% ===== OUTPUT DESCRIPTION =====
\outputDesc{}

\begin{itemize}
    \item In ra $n$ số nguyên $p_1, p_2, \dots, p_n$.
    \item $p_i$ là cha của đỉnh $i$ trong cây phân rã trọng tâm.
    \item Nếu $i$ là gốc của cây phân rã thì in $-1$.
\end{itemize}

Nếu có nhiều cây phân rã trọng tâm hợp lệ, in ra một cây bất kỳ.

% ===== SUBTASKS =====
% \noindent{\sffamily \textbf{Subtasks}}
% \begin{itemize}
%     \item \textbf{Subtask ?} ($?\%$ số điểm): Không có ràng buộc gì thêm.
% \end{itemize}

\vspace{0.5em}

\ExampleBox{1}{
7\\
1 2\\
1 3\\
2 4\\
2 5\\
3 6\\
3 7
}{
-1 1 1 2 2 3 3
}{}

\ExampleBox{2}{
10\\
1 2\\
2 3\\
3 4\\
4 5\\
3 6\\
6 7\\
6 8\\
8 9\\
8 10
}{
2 3 -1 3 4 8 6 3 8 8
}{}

\newpage
\def\ProblemCode{DPOPT}
\def\ProblemTitle{DP Optimization (D\&C / CHT)}
\makeproblem{\ProblemCode}{\ProblemTitle}

Cho dãy $n$ phần tử $a_1,a_2,\dots,a_n$. Ta định nghĩa giá trị $C(l,r)$ của đoạn con liên tiếp từ $l$ tới $r$ là bình phương tổng các phần tử trong đoạn đó, cụ thể là $C(l,r)=(a_{l} + a_{l+1} + \dots + a_r)^2$.

Hãy chia dãy thành đúng $k$ đoạn con liên tiếp sao cho mỗi phần tử thuộc đúng một đoạn con và tổng giá trị tất cả các đoạn đó là nhỏ nhất.

% ===== INPUT DESCRIPTION =====
\inputDesc{}

\begin{itemize}
    \item Dòng đầu tiên gồm hai số nguyên dương $n, k\ (1 \le n \le 10^{4}, 1 \le k \le \min(100,n))$.
    \item Dòng tiếp theo    gồm $n$ số nguyên $a_1,a_2,\dots,a_n\ (1 \le a_i \le 10^3)$.
\end{itemize}

% ===== OUTPUT DESCRIPTION =====
\outputDesc{}

\begin{itemize}
    \item Gồm một dòng duy nhất chứa tổng giá trị của các đoạn con nhỏ nhất tìm được.
\end{itemize}

% ===== SUBTASKS =====
% \noindent{\sffamily \textbf{Subtasks}}
% \begin{itemize}
%     \item \textbf{Subtask ?} ($?\%$ số điểm): Không có ràng buộc gì thêm.
% \end{itemize}

\vspace{0.5em}

\ExampleBox{1}{
5 3\\
1 1 1 1 1
}{
9
}{
Ta có thể chia dãy như sau: $\{1,1\ |\ 1, 1\ |\ 1\}$. \\
Tổng giá trị là $2^2+2^2+1=9$, nhỏ nhất tìm được.
}

\ExampleBox{2}{
5 3\\
8 9 5 9 7
}{
516
}{
Ta có thể chia dãy như sau: $\{8\ |\ 9,5\ |\ 9,7\}$. \\
}


\newpage
\def\ProblemCode{GEOMETRY}
\def\ProblemTitle{Hình học cơ bản}
\makeproblem{\ProblemCode}{\ProblemTitle}

Cho $n$ điểm $P_1,P_2,\dots,P_n$ trên mặt phẳng tọa độ, được đánh số từ $1$ đến $n$.  
Hãy xử lý $q$ truy vấn hình học trên các điểm này.

Các truy vấn có một trong các dạng sau:

\begin{itemize}
    \item \boxed{\texttt{ADD}\ i\ j}: In ra tọa độ của $\vec{P_i} + \vec{P_j}$.

    \item \boxed{\texttt{SCALE}\ i\ k}: In ra tọa độ của $k \cdot \vec{P_i}$.

    \item \boxed{\texttt{EQUAL}\ i\ j}: Kiểm tra hai điểm $P_i$ và $P_j$ có trùng nhau hay không.

    \item \boxed{\texttt{DOT}\ i\ j}: In ra tích vô hướng $\vec{P_i} \cdot \vec{P_j}$.

    \item \boxed{\texttt{CROSS}\ i\ j}: In ra tích có hướng $\vec{P_i} \times \vec{P_j}$.

    \item \boxed{\texttt{CCW}\ a\ b\ c}: Xác định hướng quay của ba điểm $P_a, P_b, P_c$.
    
    In ra:
    \begin{itemize}
        \item \boxed{\texttt{1}} nếu $P_a \rightarrow P_b \rightarrow P_c$ quay ngược chiều kim đồng hồ.
        \item \boxed{\texttt{-1}} nếu quay theo chiều kim đồng hồ.
        \item \boxed{\texttt{0}} nếu ba điểm thẳng hàng.
    \end{itemize}

    \item \boxed{\texttt{DIST}\ i\ j}: In ra khoảng cách giữa $P_i$ và $P_j$.

    \item \boxed{\texttt{ONSEG}\ p\ a\ b}: Kiểm tra $P_p$ có nằm trên đoạn $P_aP_b$ hay không $(P_a \ne P_b)$.

    \item \boxed{\texttt{DISTSEG}\ p\ a\ b}: In ra khoảng cách từ $P_p$ đến đoạn $P_aP_b$. $(P_a \ne P_b)$.

    \item \boxed{\texttt{PROJECT}\ p\ a\ b}: In ra hình chiếu vuông góc của $P_p$ lên đường thẳng $P_aP_b$ $(P_a \ne P_b)$.

    \item \boxed{\texttt{REFLECT}\ p\ a\ b}: In ra điểm đối xứng của $P_p$ qua đường thẳng $P_aP_b$ $(P_a \ne P_b)$.

    \item \boxed{\texttt{INTERSECT}\ a\ b\ c\ d}: Kiểm tra hai đoạn $P_aP_b$ và $P_cP_d$ có cắt nhau hay không.
    
    Nếu chúng cắt nhau tại đúng một điểm $P$, in \texttt{YES} rồi tọa độ $P$.
    Nếu không, in \texttt{NO}.
    
    \textbf{Lưu ý:} Nếu hai đoạn thẳng trùng nhau hoặc chồng lấp trên vô số điểm,
    truy vấn này cũng được xem là \texttt{NO}.
\end{itemize}

% ===== INPUT DESCRIPTION =====
\inputDesc{}

\begin{itemize}
    \item Dòng đầu tiên gồm hai số nguyên $n, q\ (1 \le n, q \le 10^5)$.
    \item $n$ dòng tiếp theo, dòng thứ $i$ chứa hai số thực $x_i, y_i\ (-10^6 \le x_i,y_i \le 10^6)$ là tọa độ của điểm $i$. Phần thập phân của mỗi tọa độ có nhiều nhất $6$ chữ số.
    \item $q$ dòng tiếp theo, mỗi dòng chứa một truy vấn theo các dạng đã mô tả ở trên.
\end{itemize}

% ===== OUTPUT DESCRIPTION =====
\outputDesc{}

\begin{itemize}
    \item Với mỗi truy vấn, in ra kết quả tương ứng trên một dòng.
\end{itemize}

Các giá trị thực sẽ được chấp nhận nếu sai số tuyệt đối không vượt quá $10^{-3}$,
tức là nếu đáp án đúng là $x$ và bạn in ra $y$ thì $|x-y| \le 10^{-3}$.

\vspace{0.5em}

\ExampleBox{1}{
5 15\\
0 0\\
2 0\\
2 2\\
0 2\\
1 3\\
ADD 2 4\\
SCALE 3 1.5\\
EQUAL 5 3\\
DOT 2 4\\
CROSS 2 4\\
CCW 3 4 5\\
DIST 1 3\\
ONSEG 5 3 4\\
DISTSEG 5 1 4\\
DISTSEG 3 1 4\\
PROJECT 5 1 4\\
REFLECT 5 1 4\\
REFLECT 5 1 3\\
INTERSECT 1 3 2 4\\
INTERSECT 4 5 1 2
}{
2.000000000000 2.000000000000\\
3.000000000000 3.000000000000\\
NO\\
0.000000000000\\
4.000000000000\\
-1\\
2.828427124746\\
NO\\
1.414213562373\\
2.000000000000\\
0.000000000000 3.000000000000\\
-1.000000000000 3.000000000000\\
3.000000000000 1.000000000000\\
YES 1.000000000000 1.000000000000\\
NO
}{}

% TẠO BÀI TẬP MỚI:
% \newpage
% \input{<tên file bài>}

%==========================

\label{LastProblem}
\label{LastPage}

\vfill
\begin{center}
    \textbf{*** HẾT ***}
\end{center}
\vspace{3em}

\end{document}